\begin{problem}[33届物理竞赛决赛]{69}
    六、Hyperloop 是一款利用胶囊状的运输车在水平管道中的快速运动来实现超高速运输的系统（见图 a）。它采用了“气垫”技术和“直线电动机”原理。

    “气垫”技术是将内部高压气体从水平放置的运输车下半部的细孔快速喷出（见图 b），以至于整个运输车被托离管壁非常小的距离，从而可忽略摩擦。运输车横截面是半径为 $R$ 的圆，运输车下半部壁上均匀分布有沿径向的大量细孔，单位面积内细孔个数为 $n$ （$n >> 1$ ），单个细孔面积为 $s$ 。运输车长度为 $l$ ，质量为 $M$ 。气体的流动可认为遵从伯努利方程，且温度不变，细孔出口处气体的压强为较低的环境压强 $P_{\mathrm{low}}$ 。

    如图 c，在水平管道中固定有两条平行的水平光滑供电导轨（粗实线），运输车上固定有与导轨垂直的两根导线（细实线）；导轨横截面为圆形，半径为 $r_{d}$ ，电阻率为 $\rho_{d}$ ，两导轨轴线间距为 $2(D+r_{d})$ ；两根导线的粗细可忽略，间距为 $D$；每根导线电阻是长度为 $D$ 的导轨电阻的2倍。两导线和导轨轴线均处于同一水平面内。导轨、导线电接触良好，且所有接触电阻均可忽略。
    \begin{subproblem}
        求内部高压气体压强 $P$ 为多大时运输车才刚好能被托起？
    \end{subproblem}
    \begin{subproblem}
        如图 d 所示，当运输车进站时，运输车以速度 $v_{0}$ 沿水平光滑导轨滑进匀强磁场区域，磁场边界与导线平行，磁感应强度大小为 $B$，方向垂直于导轨-导线平面向下。当两根导线全都进入磁场后，求运输车滑动速度的大小。
    \end{subproblem}
    \begin{subproblem}
        当运输车静止在水平光滑导轨上准备离站时，在导线2后间距为 $D$ 处接通固定在导轨上电动势为 $\mathcal{E}$ 的恒压直流电源（如图 e 所示）。设电源体积及其所连导线的电阻可忽略，求刚接通电源时运输车的加速度大小。

        已知某恒流闭合回路中的一圆柱形直导线段，电流沿横截面均匀分布，如图 f 所示，其在空间中距导线轴线距离为 $r_{0}$ 的某点产生的磁感应强度方向垂直于此点和导线轴线构成的平面，大小可用下式近似计算
        \begin{equation*}
            B = \frac{\mu_{0} I}{4\pi r_{0}} (\cos\theta_{1} + \cos\theta_{2})
        \end{equation*}
        其中 $I$ 为电流，$\theta_{1}$ 、$\theta_{2}$ 是此点与导线段轴线两端连线与导线轴线的夹角。

        可能会用到的积分公式：
        \begin{equation*}
            \int_{a}^{b} \frac{1}{x} \frac{c}{\sqrt{x^{2} + c^{2}}} \mathrm{d}x = \ln\left( \frac{b}{a} \frac{c + \sqrt{c^{2} + a^{2}}}{c + \sqrt{c^{2} + b^{2}}} \right), \quad \text{其中 } a、b、c \text{ 均为正数。}
        \end{equation*}
    \end{subproblem}
\end{problem}